\section{Spectral rigidity}
\label{sec:rigidity}

We now prove \Cref{thm:rigidity}. Throughout this Section, assume its
hypotheses and retain the modular notation and the symmetric mean from
\Cref{mod:average}. In particular, $M_\phi=\C1$. We use no ergodicity
assumption on $b$.

\subsection{From bounded recovery to an interval estimate}

When $a\xi=\xi a^\flat$, the identity $wa\xi=R(a^\flat)w\xi$ for
$w\in M$ writes the intertwining defect, at fixed $s$, as a bounded
operator applied to $w\xi$.

\begin{lemma}[Estimate at a spectral point]
\label{spec:point}
Let $a,a^\flat\in M$ satisfy $a\xi=\xi a^\flat$. For $s\in\R$, set
\[
 T_s=R(a^\flat)-L(b_s(a))\in\B(H).
\]
If $h_n\in H$ are unit vectors with $D$-spectral supports contained in
$[s-d_n,s+d_n]$, where $d_n>0$ and $d_n\to0$, then
\begin{equation}
 \norm{t\mapsto T_sU_th_n}_{\mathrm{av}}\longrightarrow0.
 \label{eq:point-estimate}
\end{equation}
\end{lemma}

\begin{proof}
If the conclusion fails, \Cref{rec:bounded}, applied to the fixed operator
$T_s$, gives a uniformly bounded sequence $v_n\in M$ such that the
$D$-spectral supports of $v_n\xi$ shrink to $s$ and
\[
 \norm{T_sv_n\xi}\geq\eta>0
 \qquad\text{for every }n.
\]
By \Cref{mod:predual},
\[
 \norm{\phi v_n-e^s v_n\phi}_{M_*}\longrightarrow0.
\]
The intertwining hypothesis \eqref{eq:flow-property} therefore gives
$(v_na-b_s(a)v_n)\xi\to0$. Since left and right multiplication commute,
\[
 T_sv_n\xi
   =v_n\xi a^\flat-b_s(a)v_n\xi
   =(v_na-b_s(a)v_n)\xi,
\]
a contradiction.
\end{proof}

The point estimate supplies no uniform rate in $s$. To pass to an interval,
we use the commutation of the averaged quadratic form with the spectral
projections of $D$.

\begin{lemma}[Estimate on an interval]
\label{spec:interval}
Let $a,a^\flat\in M$ satisfy $a\xi=\xi a^\flat$. For every nonempty compact
interval $I\subset\R$, every $c\in\C$, and every $h\in E_D(I)H$,
\begin{equation}
 \norm{t\mapsto (R(a^\flat)-cL(a))U_th}_{\mathrm{av}}
 \leq \norm{h}\sup_{s\in I}\norm{(b_s(a)-ca)\xi}.
 \label{eq:interval-estimate}
\end{equation}
\end{lemma}

\begin{proof}
Put
\[
 S=R(a^\flat)-cL(a),\qquad
 C=P(S^*S),\qquad
 B=\sup_{s\in I}\norm{(b_s(a)-ca)\xi}.
\]
The number $B$ is finite, and \Cref{mod:average} gives
\[
 \norm{t\mapsto SU_th}_{\mathrm{av}}^2=\langle Ch,h\rangle.
\]
Suppose that the asserted estimate fails. By homogeneity there is a unit
vector $h\in E_D(I)H$ with
\[
 \langle Ch,h\rangle=B^2+\varepsilon
 \qquad(\varepsilon>0).
\]
For each positive integer $n$, partition $I$ into finitely many disjoint
Borel intervals $J$ of diameter at most $1/n$. Since $C$ commutes with the
spectral projections of $D$,
\[
 \langle Ch,h\rangle
   =\sum_J\langle C E_D(J)h,E_D(J)h\rangle,
 \qquad
 \sum_J\norm{E_D(J)h}^2=1.
\]
Consequently, one nonzero spectral piece, after normalization, gives a unit
vector $h_n$ supported in $\overline{J_n}$ such that
\begin{equation}
 \langle Ch_n,h_n\rangle\geq B^2+\varepsilon.
 \label{eq:partition-gap}
\end{equation}
By compactness of $I$, after passing to a subsequence these supports lie in
intervals shrinking to a single point $s\in I$. This also covers a singleton
interval $I$.

Write $z=b_s(a)-ca$, so that $S=T_s+L(z)$. By \Cref{spec:point},
$\norm{t\mapsto T_sU_th_n}_{\mathrm{av}}\to0$. On the other hand,
\Cref{mod:average} gives the exact identity
\[
 \norm{t\mapsto L(z)U_th_n}_{\mathrm{av}}^2
 =\langle P(L(z^*z))h_n,h_n\rangle
 =\phi(z^*z)\leq B^2.
\]
The triangle inequality now yields
$\limsup_n\norm{t\mapsto SU_th_n}_{\mathrm{av}}\leq B$, contrary to
\eqref{eq:partition-gap}.
\end{proof}

\subsection{The joint spectral pairing}

State preservation gives a unitary implementation
\[
 V_s(a\xi)=b_s(a)\xi\qquad(a\in M).
\]
It is strongly continuous by the $u$-continuity of $b$: for $a\in M$,
\[
 \norm{V_s(a\xi)-a\xi}^2
 =2\phi(a^*a)-2\operatorname{Re}\phi(a^*b_s(a))\longrightarrow0
 \qquad(s\to0),
\]
and density extends this continuity to $H$. Since $b_s$ preserves $\phi$,
it commutes with the modular automorphism group. Thus $V_sU_t=U_tV_s$
for all $s,t$.
Write $V_s=e^{isQ}$. The self-adjoint operators $D$ and $Q$ strongly commute;
let $E$ be their joint spectral resolution and let $\mathcal S\subset\R^2$
be its closed support.

The joint spectral measure is concentrated on $\mathcal S$. Indeed, its
complement is a union of open sets with zero spectral projection, and the
countable base of $\R^2$ reduces this to a countable union. Hence
$E(\mathcal S)=1$, without any separability assumption on $H$.

\begin{lemma}[Joint spectral pairing]
\label{spec:pairing}
For every $(r,p),(q,\ell)\in\mathcal S$,
\begin{equation}
 \exp\bigl(i(r\ell+qp)\bigr)=1.
 \label{eq:spectral-pairing}
\end{equation}
\end{lemma}

\begin{proof}
Fix $\delta>0$ and write $I_v=[v-\delta,v+\delta]$. We first choose nonzero
$x,a\in M$ such that
\[
 x\xi\in E(I_r\times I_p)H,
 \qquad
 a\xi\in E(I_q\times I_\ell)H.
\]
Here is the justification for this bounded localization. Choose nonnegative
smooth functions $f,g$ supported in the interiors of $I_r,I_p$ and equal
to one on smaller neighborhoods of $r,p$. The spectral projection of the
resulting smaller rectangle is nonzero by the definition of $\mathcal S$.
Consequently, $f(D)g(Q)\neq0$. Density of $M\xi$ supplies $y\in M$ with
$f(D)g(Q)y\xi\neq0$. Fourier integration along $b$ realizes $g(Q)y\xi$ as
a vector in $M\xi$: if $g(u)=\int_{\R}\widehat g(t)e^{itu}\,dt$, use the
bounded weak integral $\int_{\R}\widehat g(t)b_t(y)\,dt$. The subsequent
filter $f(D)$ also preserves $M\xi$ by \Cref{mod:localization}. Their
composition therefore gives the required $x$. The construction of $a$ is
identical. No bound uniform in $\delta$ on the resulting operators is
needed.

Both vectors have compact modular spectral support. By
\Cref{mod:localization}, choose $x^\flat,a^\flat\in M$ with
\[
 x\xi=\xi x^\flat,\qquad a\xi=\xi a^\flat.
\]
Set
\[
 c=e^{ir\ell},\qquad d=e^{iqp},\qquad
 H_0=\norm{a\xi}\norm{x\xi}>0,
\]
and consider the bounded continuous vector functions
\[
 F(t)=\sigma_t(x)a\xi,\qquad G(t)=a\sigma_t(x)\xi.
\]
Since
\[
 F(t)-cG(t)=(R(a^\flat)-cL(a))U_t(x\xi),
\]
\Cref{spec:interval} and the $Q$-spectral support of $a\xi$ give
\begin{align}
 \norm{F-cG}_{\mathrm{av}}
 &\leq\norm{x\xi}\sup_{s\in I_r}\norm{(b_s(a)-ca)\xi}\notag\\
 &\leq H_0\sup_{\substack{s\in I_r\\u\in I_\ell}}
                \abs{e^{isu}-c}.
 \label{eq:first-product-estimate}
\end{align}
Interchanging $a$ and $x$ in the same argument yields
\[
 \norm{t\mapsto\sigma_t(a)x\xi-dx\sigma_t(a)\xi}_{\mathrm{av}}
 \leq H_0\sup_{\substack{s\in I_q\\u\in I_p}}\abs{e^{isu}-d}.
\]
The pointwise identity
\begin{equation}
 U_t\bigl(\sigma_{-t}(a)x\xi-dx\sigma_{-t}(a)\xi\bigr)
       =G(t)-dF(t)
 \label{eq:time-reversal}
\end{equation}
and reflection invariance of the mean identify the seminorm on the left
with $\norm{G-dF}_{\mathrm{av}}$. In particular, the phase is still $d$:
the unitary $U_t$ is complex linear. We have proved
\begin{equation}
 \norm{G-dF}_{\mathrm{av}}
 \leq H_0\sup_{\substack{s\in I_q\\u\in I_p}}\abs{e^{isu}-d}.
 \label{eq:second-product-estimate}
\end{equation}

The averaged norm of $F$ has the same factor $H_0$ as the two error bounds.
By reflection invariance and
\Cref{mod:average},
\begin{align}
 \norm{F}_{\mathrm{av}}^2
 &=m_t\langle L(\sigma_t(x^*x))a\xi,a\xi\rangle\notag\\
 &=\langle P(L(x^*x))a\xi,a\xi\rangle
   =\phi(x^*x)\phi(a^*a)=H_0^2.
 \label{eq:product-norm}
\end{align}
Using $\abs{c}=1$ and
\[
 (1-cd)F=(F-cG)+c(G-dF),
\]
combine \eqref{eq:first-product-estimate} and
\eqref{eq:second-product-estimate}, and divide by $H_0$, to obtain
\[
 \abs{1-cd}
 \leq
 \sup_{\substack{s\in I_r\\u\in I_\ell}}\abs{e^{isu}-e^{ir\ell}}
 +\sup_{\substack{s\in I_q\\u\in I_p}}\abs{e^{isu}-e^{iqp}}.
\]
The right side tends to zero as $\delta\to0$. For example, the first
supremum is at most $\delta(\abs{r}+\abs{\ell})+\delta^2$, by
$\abs{e^{iv}-e^{iw}}\leq\abs{v-w}$ for real $v,w$. The second has the
analogous bound. Since $c$ and $d$ depend only on the original support
points, this proves \eqref{eq:spectral-pairing}.
\end{proof}

\subsection{Conclusion of the proof}

We first record the restriction that a scalar modular centralizer imposes
on modular eigenvalues. This is the weak-mixing argument of
Marrakchi and Vaes \cite[Lemma~2.1]{MarrakchiVaes2024}; we include the
needed conclusion and its proof in our spectral notation. The bounded
operator extraction is useful because a spectral projection on $H$ need
not itself preserve $M\xi$.

\begin{lemma}[Modular eigenvalues]
\label{spec:eigenvalues}
If $M_\phi=\C1$, then $D$ has no nonzero eigenvalue. If, in addition, $D$
has pure point spectrum, then $M=\C1$.
\end{lemma}

\begin{proof}
Suppose $E_D(\{r\})\neq0$. Choose $x\in M$ such that
$\eta=E_D(\{r\})x\xi\neq0$, using density of $M\xi$. Let
\[
 y_k=\int_{\R}e^{-irt}\sigma_t(x)\,d\mu_k(t),
\]
where $\mu_k$ is uniform probability measure on $[-k,k]$. These bounded
weak integrals satisfy $\norm{y_k}\leq\norm{x}$. The Hilbert-space mean
ergodic theorem for $t\mapsto e^{-irt}U_t$ gives
\[
 y_k\xi\longrightarrow\eta\qquad\text{in }H.
\]
Take a weak operator convergent subnet of $(y_k)$, with limit $y\in M$.
Its vector limit is $y\xi=\eta\neq0$. Therefore, for every $t\in\R$,
\[
 \sigma_t(y)\xi=U_ty\xi=e^{irt}y\xi.
\]
Since $\xi$ is separating, $\sigma_t(y)=e^{irt}y$.

It follows that $y^*y,yy^*\in M_\phi=\C1$. Their scalar values coincide
because $\norm{y^*y}=\norm{yy^*}$, and are strictly positive. Thus $y$ is
a nonzero scalar multiple of a unitary, and
$\phi(y^*y)=\phi(yy^*)>0$. The modular identity
\eqref{eq:modular-identities} now gives
\[
 \phi(yy^*)
   =\norm{\xi y}^2
   =\norm{e^{D/2}y\xi}^2
   =e^r\norm{y\xi}^2
   =e^r\phi(y^*y).
\]
Hence $r=0$. If $D$ has pure point spectrum, all of $H$ is consequently
its zero eigenspace, so $D=0$. Then the modular action is trivial and
$M=M_\phi=\C1$.
\end{proof}

\begin{proof}[Proof of \Cref{thm:rigidity}]
By \Cref{spec:pairing}, the symmetric bilinear form
\[
 \mathfrak b((r,p),(q,\ell))=r\ell+qp
\]
takes values in $2\pi\Z$ on every pair of points of $\mathcal S$. We claim
that $\mathcal S$ is countable or is contained in a coordinate axis.

If $\mathcal S$ contains linearly independent vectors $v_1,v_2$, the map
\[
 z\longmapsto\bigl(\mathfrak b(v_1,z),\mathfrak b(v_2,z)\bigr)
\]
is an invertible real linear map on $\R^2$, since the form is nondegenerate
and $v_1,v_2$ are independent. It sends $\mathcal S$ into the
countable set $(2\pi\Z)^2$, proving countability. Otherwise $\mathcal S$
is contained in a line through the origin. If this line is not a coordinate
axis, write its points as $t(\alpha,\beta)$ with $\alpha\beta\neq0$.
Pairing each point with itself gives
\[
 \alpha\beta t^2\in\pi\Z.
\]
There are at most two real values of $t$ for each integer, again proving
countability. The zero-dimensional case is contained in both axes. This
proves the claim.

If $\mathcal S$ is countable, its spectral projections at singletons sum
strongly to $1$, because $E(\mathcal S)=1$. Each such projection has range
in an eigenspace of $D$, so $D$ has pure point spectrum. By
\Cref{spec:eigenvalues}, $M=\C1$, and $b$ is trivial.

If instead $\mathcal S\subset\{0\}\times\R$, then $D=0$, again giving
$M=M_\phi=\C1$ and triviality of $b$. In the remaining case
$\mathcal S\subset\R\times\{0\}$, we have $Q=0$. Thus $V_s=1$ for every
$s$, and
\[
 b_s(a)\xi=V_s(a\xi)=a\xi\qquad(a\in M).
\]
Separation by $\xi$ gives $b_s(a)=a$. In all cases $b$ is the trivial action,
as asserted.
\end{proof}
